% ============================================
% Supplementary Note 2: Assembly History and Parental-Affiliation-Correlated Species Fates
% ============================================
\subsection*{Supplementary Note 2: Assembly history and parental-affiliation-correlated species fates}
\label{sec:suppl:note2}

\paragraph{Assembly history.}
To test whether assembly history contributes to community-level selection~\citep{Gilpin1994,Debray2022,Tikhonov2016,DiazColunga2022}, we compared coalescence outcomes to a ``direct assembly'' control. In coalescence, two parental communities are first assembled separately from the species pool and then mixed. In direct assembly, all species from both pools are mixed simultaneously without prior assembly history.

In simulations, species compete during assembly and some go extinct. Those that survive to coexistence tend to have weaker mutual interactions on average. We quantified this assembly effect by measuring the average pairwise interaction coefficient $\langle \alpha_{ij} \rangle$ among species at different stages (Supplementary Fig.~32) and by visualizing pre- and post-assembly interaction submatrices (Supplementary Fig.~3). Across all tested interaction-strength parameter values ($\mu = 0$ to $1.2$), post-assembly communities exhibited substantially reduced average pairwise interaction coefficients compared to the pre-assembly pool. This reduction becomes more pronounced at higher $\mu$, where stronger competition drives more species to extinction, leaving only those with weaker mutual interactions. The assembly process thus creates internal coherence within each parental community: species that survived together share compatible (less competitive) interactions.

We hypothesized that this shared history would increase Dominance probability compared to direct assembly. Simulations supported this prediction (Extended Data Fig.~8): across interaction-strength parameter values $\mu = 0.2$ to $1.0$, coalescence consistently produced higher Dominance fractions than direct assembly. At $\mu = 0.2$, coalescence yielded 37\% Dominance versus 12\% for direct assembly. At $\mu = 0.6$, the values were 59\% versus 20\%. At $\mu = 1.0$, they were 65\% versus 22\%. Correspondingly, direct assembly showed higher Restructuring rates, as expected when species from different pools have not been pre-filtered for compatibility.

Experimental results corroborate these findings (Supplementary Fig.~23). Under Base and Nutr$+$ conditions, coalescence showed elevated Dominance rates compared to direct assembly, consistent with simulation predictions. The Nutr$-$ condition showed a weaker coalescence-versus-direct-assembly contrast, consistent with weaker invasion resistance in nutrient-limited environments (see Supplementary Note~4).

\paragraph{Pairwise selection correlation.}
To quantify whether species from the same parental community exhibit correlated selection during coalescence, we developed a pairwise fate-concordance metric. For each coalescence event, species were assigned a parental-affiliation label (\rev{the parental community in which the species had higher pre-coalescence relative abundance, if present in both}). For each species pair, we evaluated whether their presence/absence patterns in the offspring showed shared fates (both present or both absent) or different fates (one survives and the other goes extinct).

The shared-fate rate was converted to a correlation-style concordance metric: $\rho = 2 \times \text{shared-fate rate} - 1$, yielding values from $-1$ (all pairs have different fates) to $+1$ (all pairs share fate). We computed $\rho_{\text{same}}$ for within-parental-affiliation pairs and $\rho_{\text{cross}}$ for cross-affiliation pairs, then averaged across coalescence events. To test significance, we generated null distributions by shuffling parental-affiliation labels (1,000 permutations) and computed permutation $p$-values for $\Delta = \bar{\rho}_{\text{same}} - \bar{\rho}_{\text{cross}}$.

Positive $\bar{\rho}_{\text{same}}$ indicates that species with the same parental-affiliation label share fates (survive or go extinct together). Negative $\bar{\rho}_{\text{cross}}$ indicates that cross-affiliation species have opposite fates. In simulations (Extended Data Fig.~5a--d), same-parent pairs remained positively correlated, whereas cross-community pairs became increasingly negative as the interaction-strength parameter increased. At low interaction-strength parameter values ($\mu = 0.3$), both within-community and cross-community pairs show positive correlations with only a small difference ($\Delta = 0.05$; Extended Data Fig.~5a). As the interaction-strength parameter increases, the difference becomes pronounced. At $\mu = 0.6$, within-community pairs show positive correlation while cross-community pairs become negative ($\Delta = 0.58$; Extended Data Fig.~5b). At $\mu = 0.8$, the separation is larger ($\Delta = 0.94$; Extended Data Fig.~5c). In experiments (Extended Data Fig.~6), the same-versus-cross separation was strongest and statistically significant in Nutr$+$, not statistically significant in Nutr$-$, and intermediate in Base. This supports the interpretation that assembly history can structure pair-level species fates during coalescence, providing outcome-level evidence for community-level selection~\citep{Tikhonov2016,DiazColunga2022,Mansour2018,Rillig2015}.

\rev{This metric can also be interpreted through the invasion-fitness framework. Related pairwise invasion-resistance assays are summarized in Supplementary Note~4. In a gLV model, invasion fitness is the per-capita growth rate of a rare species introduced into a resident community background. Positive invasion fitness predicts establishment, whereas negative invasion fitness predicts exclusion~\citep{Hu2025,Kurkjian2021}. The 95:5 pairwise invasion assays provide an empirical two-species approximation of this idea, while the pairwise fate-concordance metric asks whether analogous persistence outcomes become correlated across many species from the same assembled parental community during coalescence. Consistent with this interpretation, an auxiliary gLV analysis showed that excess same-parent shared fate, measured relative to a parental-affiliation-shuffled null, increased with the interaction-strength parameter $\mu$ across the full simulated sweep ($\mu = 0.05$--$1.20$; Pearson $r = 0.870$, $p = 3.2 \times 10^{-8}$; Supplementary Fig.~36). This auxiliary within-parent metric measures the same parental-affiliation-correlated shared-fate signal as $\Delta=\bar{\rho}_{\text{same}}-\bar{\rho}_{\text{cross}}$, but compares same-parent pairs to a parental-affiliation-shuffled null rather than explicitly subtracting cross-community pairs.}
